\documentclass[11pt,fontset=none]{ctexbook}

\usepackage[paper=b5paper,twoside=true]{geometry}
\usepackage{graphicx}
\usepackage[affil-it]{authblk}
\usepackage{indentfirst}
\usepackage[shortlabels]{enumitem}
\setlist{nosep}
\newcommand{\rmparen}[1]{{\rm (}#1{\rm )}}

\usepackage{amsmath}
\usepackage{amsthm}
\usepackage{amssymb}
\usepackage[all]{xy}
\usepackage{pgfplots}
\numberwithin{equation}{section}
\allowdisplaybreaks
\pgfplotsset{compat=1.17}
\tikzset{every picture/.style={line width=0.75pt}} %set default line width to 0.75pt

\DeclareMathOperator*{\grad}{grad}
\DeclareMathOperator*{\Hom}{Hom}
\DeclareMathOperator*{\Obj}{Obj} % 我承诺这篇笔记中不会有范畴论
\DeclareMathOperator*{\codim}{codim}
\DeclareMathOperator{\Aut}{Aut}
\DeclareMathOperator{\End}{End}
\DeclareMathOperator*{\Span}{Span}

\renewcommand{\proofname}{{\bf 证明}}
\theoremstyle{definition}
\newtheorem{defn}{定义}[section]
\newtheorem{sym}[defn]{记号}
\newtheorem{eg}[defn]{例}
\theoremstyle{plain}
\newtheorem{thm}[defn]{定理}
\newtheorem{lem}[defn]{引理}
\newtheorem{col}[defn]{推论}
\newtheorem{prop}[defn]{命题}
\newtheorem*{pro}{问题}
\theoremstyle{remark}
\newtheorem{rem}[defn]{评注}
\newtheorem{ex}[defn]{习题}

\setmainfont{TeX Gyre Termes}
\setsansfont{TeX Gyre Heros}
\setmonofont{TeX Gyre Bonum}
\newfontfamily\greek{FreeSerif}
\usepackage{unicode-math}
\setmathfont{TeX Gyre Termes Math}
\setmathfont{STIX Two Math}[range={"2211,"222B,"1D70B}]
\setmathfont{TeX Gyre Pagella Math}[range={bb,"1D51B}]
\setCJKmainfont[BoldFont={FZCSJ.otf},
    ItalicFont={FZKTJ.otf},
    BoldItalicFont={FZCKJ.otf}]{FZSSJ.otf}
\setCJKsansfont{FZHTJ.otf}
\setCJKmonofont{FZFSJ.otf}

\usepackage{titlesec}
\titleformat{\chapter}[frame]{\Huge\bfseries}{第\,\thechapter\,章}{1em}{}{}
\titleformat*{\subsection}{\Large\itshape}

\usepackage[style=caspervector,sorting=ecnyt,utf8,seconds=true]{biblatex}
\addbibresource{biblio.bib}

\usepackage{xcolor}
\usepackage[colorlinks=true,linkcolor=blue,citecolor=green,urlcolor=red]{hyperref}

\title{曲面曲线论笔记}
\author{魔法少女Alkali}
\affil{北京师范大学数学科学学院}
\date{2022 年 5 月}

\begin{document}
    \maketitle
    \vspace*{13cm}
    \thispagestyle{empty}
    \begin{figure}[ht]
        \includegraphics[width=0.4\textwidth]{by-sa.pdf}\hspace{0.05\textwidth}\begin{minipage}[b]{0.55\textwidth}
        \small
        本作品采用知识共享署名-相同方式共享 4.0 国际 (CC BY-SA 4.0) 协议进行共享.
        您可以访问\url{https://creativecommons.org/licenses/by-sa/4.0/}查看该协议.
        \end{minipage}
    \end{figure}

    \frontmatter
    \chapter{前言}
还没想好写什么.

    \tableofcontents

    \mainmatter
    \setcounter{chapter}{-1}
\chapter{准备}
\section{几何学是什么?}
现代几何学源于古希腊.
在古希腊语中, ``几何学''一词为γεωμετρία \textit{(geōmetría)}, 意为测量大地.
这反映了早期几何学主要是对长度, 面积, 体积, 角度的经验性原理的收集, 主要用于满足实用性用途.
直至今日, 初步的几何学教学仍然是从对几何体的大小的直观认识开始的.
因此几何学的一个经典要件就是\emph{度量}.

Euclid所著的\textit{Elements} (汉语中通常称作《几何原本》) 是古希腊几何的代表.
他在\textit{Elements}的开篇引入了这样的一条公理:
\begin{quotation}
    {\bfseries 公理}4.\ 彼此能重合的物体是全等的.
\end{quotation}
然后第一个引用了这条公理的命题是
\begin{quotation}
    {\bfseries 命题}4.\ 如果两个三角形中, 一个的两边分别等于另一个的两边, 而且这些线段所夹的角相等.
    那么它们的底边等于底边, 这样其余的角也等于相应的角, 即那些等边所对应的角.\\
    \begin{flushright}
        (译文引自\parencite{Euclid_Elem})
    \end{flushright}
\end{quotation}
在这最原始的直觉中, ``重合''蕴含了运动的概念, 而边角的相等则蕴含了不变量的概念.
因此, 几何学的另一个经典要件就是\emph{变换与不变量}.

古希腊的几何学主要研究直线与圆锥曲线.
到了微积分发明之后, 数学家可以使用微积分的工具来研究更一般的几何体了.
Leibniz通过密切圆引入了曲线的曲率, Bernoulli与Euler研究了曲面的法曲率与测地线.
对``曲''的研究正式进入了几何学之中.
1827年, Gauss在论文\textit{Disquisitiones generales circa superficies curvas} (关于曲面的一般研究) 中证明了``Gauss绝妙定理'' (Theorema Egregium).
从此一种观念开始进入几何学: 我们可以研究抽象的几何体, 而不考虑它在欧式空间中的实现.
进而不久非欧几何便产生了.

回到几何学的两个要件上.
有了以上基础, 1854年Riemann写作了论文\textit{\"{U}ber die hypothesen welche der Geometrie zu Grunde liegen} (论奠定几何学基础的假设), 正式引入了Riemann度量与流形的概念 (之后的笔记中我们会详细解释这两个概念).
同时, 19世纪正在经历一个射影几何的复兴潮, 当时正在流行使用射影变换的方法研究射影几何.
于是在1893年F.\ Klein发表了对整个几何学的``总结性''综述\textit{Vergleichende Betrachtungen \"{u}ber neuere geometrische Forschungen} (关于近代几何学研究的比较考察).
文章中提到几何学的目的在于
\begin{quotation}
    给定一个流形和其上的变换群, 建立关于这个变换群的不变量理论.
    \begin{flushright}
        (\parencite{Klein_Erlangen}, 自译)
    \end{flushright}
\end{quotation}

自此, 经典几何学的舞台已经搭好.
但我们还有一个问题: \label{intro_sect}
\begin{pro}
    微分几何是什么?
\end{pro}
接下来我们开始慢慢搭微分几何的舞台.
    \chapter{曲线的局部理论}

我们在这一章会考虑具有平凡切丛的流形上的曲线.
我们使用的仍然是Frenet标架法, 并且会证明类似于 % R^3中的曲线论基本定理
的定理, 以对曲线在等距变换群下进行分类.
    
    \chapter{计算}
前面我们遇到的大量几何量实际上都是\textit{张量}.
我们打算在这一章提供一些必要的代数预备知识, 并介绍张量入门.
本章中所有的环都有单位元.

\section{代数知识}

\begin{defn}\label{def of left module}
    一个环$R$上的\textbf{左模}由一个Abel群$M$与$R$到$M$的自同构群的一个同态构成.
    换言之, 任意一个$r\in R$都诱导了一个$M\to M$的映射, 对$m\in M$, $m$在$r$诱导的映射下的像记为$rm$.
    我们要求这些映射满足如下公理:
    \begin{enumerate}
        \item 对$r\in R$, $m,n\in M$有$r(m+n)=rm+rn$;
        \item 对$r,s\in R$, $m\in M$有$(r+s)m=rm+sm$;
        \item 对$r,s\in R$, $m\in M$有$(rs)m=r(sm)$;
        \item 对任意$m\in M$有$1m=m$.
    \end{enumerate}
\end{defn}

\begin{defn}
    一个环$R$上的\textbf{右模}由一个Abel群$M$与$R$到$M$的自同构群的一个反同态构成.
    即将定义~\ref{def of left module}~中的公理3变为$(rs)m=s(rm)$.
\end{defn}

\begin{rem}
    如果$R$是交换的, 那么左模和右模是一样的, 称为\textbf{双侧模}或者直接简称为\textbf{模}.
    此外, 我们也可以把右模的像写成$mr$, 那么``数乘的结合公理''\footnote{打引号是因为我们只对向量空间说数乘, 模一般说作用.}可以写成很好看的形式$m(sr)=(ms)r$.
\end{rem}

    \appendix
    \printbibliography[title={参考文献},heading=bibintoc]
\end{document}